\documentclass[12pt, a4paper]{report}

% --------------------------

\usepackage[spanish, mexico]{babel}

\usepackage{biblatex}   % Bibliografía

\usepackage{dirtree}    %  Árboles de directorios

% --------------------------

\title{Análisis y Diseño}

\author{Montiel Ramírez Brayan}

\setlength{\parindent}{0pt}

% --------------------------

\begin{document}

\maketitle



\section{Computabilidad}

El problema de crear un chat que utilice el protocolo indicado es computable, ya que consiste en la creación de dos programas (cliente y servidor) que utilizan algoritmos con un número finito de pasos, los cuales son las 12 funciones que describe el protocolo, que a su vez se descomponen en algoritmos con un número finito de pasos también.

\section{Paradigma}

El protocolo detalla la comunicación entre un \textbf{cliente} y un \textbf{servidor} mediante \textbf{mensajes} y describe 12 funciones que el chat debe permitir. Esto es motivo suficiente para escoger el paradigma orientado a objetos, pues más adelante se detallarán las componentes que conforman a estos dos programas y se modelarán como objetos con sus propias responsabilidades.

\chapter{Análisis fundamental}

Seguiremos el modelo cliente-servidor, y usaremos \textit{sockets} para establecer la conexión entre estos dos programas, en particular, usaremos el protocolo ITC/IP; además, habrá una única instancia del servidor ejecutándose en el mismo equipo, por el alcance del proyecto.
Para que el cliente pueda establecer una conexión con el servidor, este último debe conectarse a una IP local en un puerto configurable (por defecto 1234) y ponerse en modo LISTEN.

\section{Servidor}

El servidor admite doce tipos de mensajes, los cuales indican una \textbf{acción} solicitada y un conjunto de \textbf{argumentos} relativos a ella, que entendemos como los datos asociados a dicha acción.
El servidor debe contar con una estructura o \textbf{base de datos} por lo menos, para almacenar los datos de los clientes que están usando el chat, como sus estados, salas, etc.

Analizando el flujo de las operaciones que el protocolo describe implícitamente, dado un cliente que se va a conectar con el servidor, se observan seis responsabilidades fundamentales para las componentes que este tenga:

\begin{enumerate}
	\item Conectar el socket del servidor, ponerlo en modo LISTEN y aceptar conexiones.
	\item Recibir los datos del socket extrayendo los bytes transmitidos desde el cliente.
	\item Decodificar los bytes en cadenas (los mensajes) bajo el formato UTF-8.
	\item Validar el formato del mensaje bajo el estándar JSON y la identificación del usuario (cuando aplique) y asignar una estrategia de acción-respuesta correspondiente al mensaje.
	\item Tener un algoritmo para acceder o modificar información dentro del servidor, notificar a otros usuarios, responder al cliente cuando sea necesario y recibir parámetros personalizados.
	\item Ejecutar la estrategia de la acción-respuesta asignada al mensaje.
\end{enumerate}

A los objetos que cumplan con estas responsabilidades les llamaremos:

\begin{enumerate}
	\item Conector
	\item Proovedor de datos de la conexión
	\item Traductor
	\item Intérprete
	\item Estrategia de acción-respuesta
	\item Ejecutor de estrategia
\end{enumerate}

respectivamente, todo relativo al servidor.

\section{Cliente}

\section{Estructura de directorios}

\dirtree{%
    .1 FC\_MYP\_Proyecto1/.
        .2 Docs/.
            .3 Process/.
            .3 Analysis\_Design/.
            .3 CodeDocs/.
                .4 Client/.
                .4 Server/.
        .2 Client/.
            .3 Model/.
            .3 View/.
            .3 Controller/.
        .2 Server/.
            .3 Model/.
            .3 View/.
            .3 Controller/.
}
\vspace{2mm}

\subsection{Directorio {\ttfamily FC\_MYP\_Proyecto1}}
Es la raíz del proyecto, y aunque sus subdirectorios tienen nombres claros, contiene además cualquier archivo que requiera estar en la raíz del proyecto, así como archivos de configuración como el {\ttfamily .gitignore}, archivos de configuración para la generación automatizada de documentación, etc.

\subsection{Directorio {\ttfamily Docs}}
Sus subdirectorios son muy descriptivos:
\begin{itemize}
    \item En {\ttfamily Process} se documenta el proceso de cómo se desarrolló el proyecto completo y qué flujos de trabajo se siguieron.
    \item En {\ttfamily Analysys\_Design} va este mismo documento y cualquier otro recurso que ayudara para completar las etapas de análisis y diseño.
    \item En {\ttfamily CodeDocs} va la documentación del código, generada por la herramienta \textit{Doxygen} mediante la integración continua. Tiene una sección para la documentación del código del cliente y otra para la del código del servidor.
\end{itemize}

\subsection{Directorio {\ttfamily Client}}
Contiene el ejecutable del cliente y utiliza el patrón de diseño MVC, dividiendo la conexión y comunicación con el servidor, la interfaz gráfica y el canal de comunicación entre ambos en tres capas diferentes.

\subsection{Directorio {\ttfamily Server}}
Contiene el ejecutable del servidor y al igual que con el cliente, usa el patrón MVC, dividiendo las reglas de la aplicación desde el servidor, la impresión en consola y la comunicación entre ambos en tres capas.

\end{document}